\clearpage
\VerdierPhysicalPageAnchor{0131}%
\begingroup
\small
\setlength{\emergencystretch}{1.5em}%
\setlength{\abovedisplayskip}{2pt plus 1pt minus 1pt}%
\setlength{\belowdisplayskip}{2pt plus 1pt minus 1pt}%

\VerdierSemanticAnchor{ch2-2-1-5-continuation}
\noindent pleine de $\mathcal D$ admet au plus une structure triangulée
compatible avec le foncteur d'inclusion.

\VerdierSemanticAnchor{ch2-2-1-5-criterion}
Soit $\mathcal D'$ une sous-catégorie pleine de $\mathcal D$ stable par
translations. Pour que $\mathcal D'$ admette une structure triangulée faisant
de $\mathcal D'$ une sous-catégorie triangulée pleine de $\mathcal D$, il faut
et il suffit que pour tout couple $X,Y$ d'objets de $\mathcal D'$ et tout
morphisme $f\colon X\to Y$, il existe un triangle distingué de $\mathcal D$ :
$(X,Y,Z,f,g,h)$, où $Z$ est un objet de $\mathcal D'$.

\VerdierSemanticAnchor{ch2-2-1-6}
\medskip
\noindent\textbf{2.1.6.}\quad
\VerdierSemanticAnchor{ch2-2-1-6-saturated}
Soit $\mathcal D$ une catégorie triangulée. Une sous-catégorie triangulée
pleine $\mathcal D'$ de $\mathcal D$ est dite \emph{saturée} si elle possède
la propriété suivante :

\VerdierSemanticAnchor{ch2-2-1-6-direct-summand}
\par\smallskip
\noindent\makebox[0pt][l]{(2.1.6.1)}%
\begin{minipage}[t]{\linewidth}
  \centering\itshape
  Tout facteur direct dans $\mathcal D$ d'un objet de $\mathcal D'$\\[-1pt]
  est isomorphe à un objet de $\mathcal D'$.
\end{minipage}
\par\smallskip

\VerdierSemanticAnchor{ch2-2-1-6-direct-summand-definition}
\noindent (Un facteur direct d'un objet $Y$ est un objet $X$ tel qu'il existe
un objet $Z$ et un isomorphisme $X\oplus Z\xrightarrow{\sim}Y$).

\VerdierSemanticAnchor{ch2-2-1-7}
\medskip
\noindent\textbf{2.1.7.}\quad
\VerdierSemanticAnchor{ch2-2-1-7-denominator-set}
Soit $\mathcal D'$ une sous-catégorie triangulée pleine d'une catégorie
triangulée $\mathcal D$. Désignons par $S(\mathcal D')$ l'ensemble de
morphismes de $\mathcal D$ défini de la manière suivante :

\VerdierSemanticAnchor{ch2-2-1-7-characterization}
\par\smallskip
\noindent\makebox[0pt][l]{(2.1.7.1)}%
\begin{minipage}[t]{\linewidth}
  \centering\itshape
  Un morphisme $f\colon X\to Y$ appartient à $S(\mathcal D')$ si et seulement s'il existe\\[-1pt]
  un triangle distingué de $\mathcal D$ : $(X,Y,Z,f,g,h)$, où $Z$ est un objet de $\mathcal D'$.
\end{minipage}
\par\smallskip

\VerdierSemanticAnchor{ch2-2-1-8}
\medskip
\noindent\textbf{Proposition 2.1.8.}\quad
\VerdierSemanticAnchor{ch2-2-1-8-proposition}
\emph{L'ensemble $S(\mathcal D')$ (2.1.7.1) est un système multiplicatif
(2.1.4) de $\mathcal D$. Il est saturé si et seulement si la sous-catégorie
$\mathcal D'$ est saturée.}

\VerdierSemanticAnchor{ch2-2-1-8-proof-plan}
Pour prouver la première assertion, il suffit en vertu de (2.1.3), de montrer
que l'ensemble $S(\mathcal D')$ possède les propriétés SM1), SM3), SM5) et
SM6). Nous allons montrer successivement que l'ensemble $S(\mathcal D')$
possède ces propriétés, et nous démontrerons ensuite la dernière assertion.

\VerdierSemanticAnchor{ch2-2-1-9}
\medskip
\noindent\textbf{2.1.9.}\quad
\VerdierSemanticAnchor{ch2-2-1-9-sm1}
\emph{L'ensemble $S(\mathcal D')$ possède la propriété SM1).}

\VerdierSemanticAnchor{ch2-2-1-9-identity}
Tout d'abord, il est clair que pour tout objet $X$ de $\mathcal D$, le
morphisme $\operatorname{id}_X$ est un élément de $S(\mathcal D')$. Ceci
résulte immédiatement de TRI) et du fait que les objets nuls de $\mathcal D'$
sont des objets nuls de $\mathcal D$.

\VerdierSemanticAnchor{ch2-2-1-9-composition}
Le fait que l'ensemble $S(\mathcal D')$ est stable par composition résulte
immédiatement de l'axiome TRIV) et de (1.2.4). Nous laissons au lecteur le soin
de faire les vérifications.

\VerdierSemanticAnchor{ch2-2-1-10}
\medskip
\noindent\textbf{2.1.10.}\quad
\VerdierSemanticAnchor{ch2-2-1-10-sm3}
\emph{L'ensemble $S(\mathcal D')$ possède la propriété SM3).}

\VerdierSemanticAnchor{ch2-2-1-10-reduction}
Il suffit tout d'abord de montrer que lorsqu'on a un diagramme
$X\xrightarrow{s}Y\xrightarrow{f}Z$, avec $s\in S(\mathcal D')$, $fs=0$, il
existe un diagramme $Y\xrightarrow{f}Z\xrightarrow{t}W$, avec
$t\in S(\mathcal D')$, $tf=0$. Car alors, en passant aux catégories opposées,
on obtiendra la réciproque et, par

\VerdierSemanticAnchor{ch2-2-1-10-p132-boundary}
\endgroup
